\begin{abstract}
Frontier-based navigation agents evaluate every frontier with the same static scoring function throughout an episode, regardless of whether their decisions prove correct. We propose PEC-Nav, a self-calibrating search mechanism for Embodied Question Answering that learns from its own prediction errors online. Before exploring a frontier, the agent predicts the semantic layout beyond it; after arriving, it measures the divergence between prediction and observation as a \emph{surprise signal}. High surprise broadens exploration; low surprise narrows the search toward the goal, an adaptive shift that emerges automatically without environment-specific training. PEC-Nav wraps this Predict-Explore-Correct loop around a language-driven EQA pipeline with episodic memory and follow-up QA. Because memory and reasoning modules are held fixed, any improvement in answer quality is directly attributable to search. On LMEE-Bench (GOAT-Bench Val Unseen, 345 episodes), PEC-Nav achieves 29.74\% SR and 20.88 SPL using a 4B-parameter VLM on a single consumer GPU. Disabling the correction loop degrades performance to 12.75\% SR; re-enabling it yields a 133\% relative improvement, confirming that prediction errors, not predictions, drive effective search.
\end{abstract}
